\documentclass[11pt]{article}

\usepackage[margin=1in]{geometry}
\usepackage{amsmath,amssymb,amsthm,mathtools}
\usepackage{booktabs,array,float}
\usepackage{microtype}
\usepackage[hidelinks]{hyperref}

\newtheorem{theorem}{Theorem}[section]
\newtheorem{lemma}[theorem]{Lemma}
\newtheorem{proposition}[theorem]{Proposition}
\newtheorem{corollary}[theorem]{Corollary}
\theoremstyle{remark}
\newtheorem{remark}[theorem]{Remark}

\newcommand{\D}{\mathbb D}
\newcommand{\C}{\mathbb C}
\newcommand{\e}{\mathrm e}
\newcommand{\Pol}{\operatorname{Pol}}

\title{Sendov's Conjecture in Degree Nine}
\author{Principia Math}
\date{}

\begin{document}
\maketitle

\begin{abstract}
We prove Sendov's conjecture for polynomials of degree nine.  After fixing a zero
$a\in[0,1]$, a hypothetical counterexample gives two reciprocal configurations,
one formed from the remaining zeros and one from the critical points.  Their
common first moment yields a strong localization inequality.  A symmetric
integral identity then reduces the remaining argument to positivity of explicit
polynomials.  The range $a\le 9/20$ is elementary; the interval
$9/20\le a\le 9/10$ is covered by eighteen exact Bernstein certificates; and
the boundary range $9/10\le a\le1$ is handled by a rescaled single certificate.
All finite computations are performed in exact rational arithmetic.
\end{abstract}

\section{Introduction}

Sendov's conjecture asserts that if all zeros of a complex polynomial lie in the
closed unit disk, then the closed unit disk centered at each zero contains a
zero of the derivative.  We prove the degree-nine case.

\begin{theorem}\label{thm:main}
Let
\[
 p(z)=\prod_{k=1}^{9}(z-z_k),\qquad |z_k|\le1.
\]
For every zero $a$ of $p$, there is a zero $\zeta$ of $p'$ such that
\[
 |a-\zeta|\le1.
\]
\end{theorem}

The proof is computer-assisted only at two finite positivity checks.  Both are
certified by expanding explicit rational polynomials in tensor-product
Bernstein bases and checking every coefficient by exact arithmetic.  The
analytic reduction and the certificate polynomials are given below; the
ancillary verifier contains no floating-point operations.

By rotation we may take the distinguished zero to be a real number
$a\in[0,1]$.  We argue by contradiction, assuming throughout that every zero
of $p'$ has distance strictly greater than $1$ from $a$.

\section{Two preliminary lemmas}

We use the Grace--Walsh--Szeg\H{o} coincidence theorem in its standard form:
if a symmetric multi-affine polynomial is evaluated at points in a convex
circular region, then the same value is attained on the diagonal at a point of
that region; see, for example, \cite[Proposition~1]{BrandenWagner}.

\begin{lemma}[separation from a zero]\label{lem:separation}
Let $P$ be a polynomial of degree $n\ge2$, let $\alpha$ be a zero of $P$, and
suppose that
\[
 |\alpha-\zeta|>\rho
\]
for every zero $\zeta$ of $P'$.  Then $\alpha$ is simple and every other zero
$\beta$ of $P$ satisfies
\[
 |\beta-\alpha|>2\rho\sin\frac{\pi}{n}.
\]
\end{lemma}

\begin{proof}
The simplicity of $\alpha$ is immediate.  Write the zeros of $P'$ as
$\zeta_1,\ldots,\zeta_{n-1}$ and set $w=\beta-\alpha$.  Since
$P(\beta)=P(\alpha)$,
\[
 0=\int_0^1\frac{P'(\alpha+tw)}{P'(\alpha)}\,dt
  =\int_0^1\prod_{j=1}^{n-1}(1-tu_j)\,dt,
 \qquad
 u_j=-\frac{w}{\alpha-\zeta_j}.
\]
The symmetric multi-affine polynomial
\[
 F(u_1,\ldots,u_{n-1})=\int_0^1\prod_{j=1}^{n-1}(1-tu_j)\,dt
\]
has diagonal
\[
 F(u,\ldots,u)=\frac{1-(1-u)^n}{nu},
\]
with the value $1$ at $u=0$.  Its diagonal zeros are
$1-\e^{2\pi i k/n}$, $1\le k<n$, and hence have modulus at least
$2\sin(\pi/n)$.  The coincidence theorem therefore implies that $F$ is
nonzero whenever all $|u_j|<2\sin(\pi/n)$.  If
$|w|\le2\rho\sin(\pi/n)$, then the strict critical-point separation gives
$|u_j|<2\sin(\pi/n)$ for every $j$, a contradiction.
\end{proof}

The second lemma is a one-dimensional convexity estimate.

\begin{lemma}\label{lem:sigma}
Let $0<m<M$, let $r_1,\ldots,r_N\in[m,M]$, and suppose
\[
 \prod_{k=1}^N r_k\ge C>m^N.
\]
Let $\nu$ be the least integer such that $M^\nu m^{N-\nu}\ge C$.  Then
\[
 \sum_{k=1}^N\frac1{r_k^2}
 \le
 \frac{N-\nu}{m^2}+\frac{\nu-1}{M^2}
 +\left(\frac{M^{\nu-1}m^{N-\nu}}{C}\right)^2.
\]
\end{lemma}

\begin{proof}
Put $r_k=\e^{t_k}$.  Since $t\mapsto\e^{-2t}$ is decreasing, we may first
reduce $\sum t_k$ to $\log C$.  On the resulting section of the box
$[\log m,\log M]^N$, the convex function $\sum\e^{-2t_k}$ is maximized at an
extreme point.  Thus all but at most one coordinate are endpoints.  By the
definition of $\nu$, the maximizing configuration has $N-\nu$ entries equal
to $m$, $\nu-1$ entries equal to $M$, and one entry equal to
$C/(M^{\nu-1}m^{N-\nu})$.  Substitution gives the result.
\end{proof}

\section{The reciprocal reduction}

Write
\[
 p(z)=(z-a)\prod_{k=1}^{8}(z-z_k),
 \qquad
 p'(z)=9\prod_{j=1}^{8}(z-\zeta_j),
\]
where multiplicities are retained.  A counterexample has $a$ simple.  Set
\[
 X_k=\frac1{a-z_k},\qquad Y_j=\frac1{a-\zeta_j},\qquad
 \mu=\frac18\sum_{j=1}^{8}Y_j=u+iv.
\]
Evaluation of $p'$ and $p''/p'$ at $a$ gives the two identities
\begin{equation}\label{eq:anchors}
 9\prod_{j=1}^{8}(a-\zeta_j)=\prod_{k=1}^{8}(a-z_k),
 \qquad
 \sum_{j=1}^{8}Y_j=2\sum_{k=1}^{8}X_k.
\end{equation}
In particular,
\begin{equation}\label{eq:mu}
 \mu=\frac14\sum_{k=1}^{8}X_k,
 \qquad |\mu|<1.
\end{equation}

Put
\[
 r_k=|a-z_k|,
 \qquad
 \sigma=\sum_{k=1}^{8}r_k^{-2}.
\]
By Lemma~\ref{lem:separation} and the unit-disk hypothesis,
\begin{equation}\label{eq:annulus}
 2\sin\frac{\pi}{9}<r_k\le1+a.
\end{equation}
Moreover, \eqref{eq:anchors} and the assumed critical-point separation imply
\begin{equation}\label{eq:product}
 \prod_{k=1}^{8}r_k>9.
\end{equation}
For exact arithmetic we use
\begin{equation}\label{eq:m0}
 m_0:=\frac{681}{1000}<2\sin\frac{\pi}{9}.
\end{equation}
Indeed, $T=2\sin(\pi/9)$ is the root in $(0,1)$ of
$3T-T^3=\sqrt3$, the function $3t-t^3$ is increasing on $[0,1]$, and
\[
 3-\left(3m_0-m_0^3\right)^2
 =\frac{16853534459219919}{10^{18}}>0.
\]

The unit-disk condition also controls the real part of $\mu$.  Since
$z_k=a-X_k^{-1}$,
\[
 2a\Re X_k+(1-a^2)|X_k|^2
 =1+(1-|z_k|^2)|X_k|^2\ge1.
\]
Summing and using \eqref{eq:mu}, we obtain the basic localization inequality
\begin{equation}\label{eq:localization}
 8au+(1-a^2)\sigma\ge8.
\end{equation}

We shall also need a uniform expansion about the mean $\mu$.  Put
\[
 \eta^2=1-|\mu|^2,
 \qquad
 D_j=Y_j-\mu.
\]
Then $\sum_jD_j=0$ and
\[
 \sum_{j=1}^{8}|D_j|^2
 =\sum_{j=1}^{8}|Y_j|^2-8|\mu|^2\le8\eta^2.
\]
If $e_m(D)$ denotes the $m$th elementary symmetric function of the $D_j$, then
\begin{equation}\label{eq:em-bounds}
 |e_m(D)|\le c_m\eta^m\qquad(2\le m\le8),
\end{equation}
where
\begin{equation}\label{eq:constants}
 (c_2,\ldots,c_8)=
 \left(4,\frac{264}{35},24,56,28,8,1\right),
 \qquad
 (\ell_2,\ldots,\ell_8)=(3,2,2,1,1,0,0).
\end{equation}
For completeness, Newton's identities give
$|e_2|\le4\eta^2$, $|e_3|\le16\sqrt2\eta^3/3<264\eta^3/35$, and
$|e_4|\le24\eta^4$.  For $m\ge5$, Maclaurin's inequality applied to
$|D_1|,\ldots,|D_8|$, followed by Cauchy--Schwarz, gives
$|e_m(D)|\le\binom8m\eta^m$.

Finally, define
\begin{equation}\label{eq:I}
 G(t)=\prod_{j=1}^{8}(1-tY_j)=\frac{p'(a-t)}{p'(a)},
 \qquad
 I=\int_0^aG(t)\,dt.
\end{equation}
Then
\begin{equation}\label{eq:I-upper}
 I=-\frac{p(0)}{p'(a)},
 \qquad
 |I|<\frac a9,
\end{equation}
because $|p(0)|\le a$ and $|p'(a)|>9$.  The remainder of the proof derives
the opposite inequality.

\section{The range \texorpdfstring{$0\le a\le9/20$}{0 <= a <= 9/20}}

The case $a=0$ is already impossible: \eqref{eq:product} contradicts
$r_k=|z_k|\le1$.  Suppose $0<a\le9/20$.  Apply
Lemma~\ref{lem:sigma} with $N=8$, $C=9$, $m=m_0$, and $M=29/20$.  Here
$\nu=7$, and exact arithmetic gives
\begin{equation}\label{eq:small-sigma}
 \sigma\le
 \frac1{m_0^2}+\frac6{(29/20)^2}
 +\left(\frac{(29/20)^6m_0}{9}\right)^2
 <\frac{1101}{200}=:S_0.
\end{equation}
By \eqref{eq:localization},
\[
 u>\frac{8-S_0(1-a^2)}{8a}.
\]
The right-hand side exceeds $1$ precisely when
\[
 S_0a^2-8a+8-S_0>0.
\]
This quadratic is decreasing on $[0,9/20]$, and at $a=9/20$ it equals
$781/80000$.  Thus $u>1$, contradicting $|\mu|<1$.

\section{The range \texorpdfstring{$9/20\le a\le9/10$}{9/20 <= a <= 9/10}}

We first isolate the analytic statement certified on each subinterval.  For
$S>4$ set
\[
 L_S(a)=\frac{8-S(1-a^2)}{8a},
\]
and, for $0\le v\le1$, define
\begin{align}
 q_S(a,\eta;v)
 &=1-\frac{8-S+Sa^2}{4}v+a^2(1-\eta^2)v^2,\label{eq:qS}\\
 h_S(a,\eta)
 &=\left(\frac S4-1\right)(1-a^2)-a^2\eta^2.\label{eq:hS}
\end{align}
Using the constants in \eqref{eq:constants}, put
\begin{equation}\label{eq:ES}
 E_S(a,\eta)=
 \sum_{m=2}^{8}c_m a^{m+1}\eta^m
 \int_0^1 v^m q_S(a,\eta;v)^{\ell_m}\,dv
\end{equation}
and
\begin{equation}\label{eq:PS}
 P_{S,\lambda}(a,\eta)=
 \frac{1-a}{9}+\frac{\eta^2}{18}
 -\frac{h_S(a,\eta)^4}{9\lambda}-E_S(a,\eta).
\end{equation}
Since the exponents $\ell_m$ are nonnegative integers,
$P_{S,\lambda}$ is a polynomial with rational coefficients whenever
$S$ and $\lambda$ are rational.

\begin{proposition}[interval certificate]\label{prop:interval}
Let $0<\alpha\le\beta<1$, and let $S,\lambda,Y$ be positive numbers such that,
for every $a\in[\alpha,\beta]$,
\begin{equation}\label{eq:interval-hyp}
 \sigma\le S,
 \qquad L_S(a)\ge\lambda,
 \qquad Y^2\ge1-\lambda^2,
 \qquad 2\lambda\ge\beta,
 \qquad
 0\le\left(\frac S4-1\right)(1-\alpha^2)\le1.
\end{equation}
If
\[
 P_{S,\lambda}(a,\eta)>0
 \qquad
 (\alpha\le a\le\beta,\ 0\le\eta\le Y),
\]
then no counterexample exists with $a\in[\alpha,\beta]$.
\end{proposition}

\begin{proof}
By \eqref{eq:localization}, $u\ge L_S(a)\ge\lambda$, and hence
$0\le\eta^2=1-|\mu|^2\le1-\lambda^2\le Y^2$.
Let
\[
 J=\int_0^a(1-t\mu)^8\,dt.
\]
For $0\le t\le a$, we have
\[
 |1-t\mu|^2=1-2tu+(1-\eta^2)t^2
 \le1-2tL_S(a)+(1-\eta^2)t^2.
\]
After $t=av$, the expression on the right is $q_S(a,\eta;v)$.  It is
nonnegative, and
\[
 q_S(a,\eta;v)-1
 \le av(-2\lambda+\beta)\le0.
\]
Thus $0\le q_S\le1$.  Expanding
\[
 G(t)=(1-t\mu)^8+
 \sum_{m=2}^{8}(-t)^m e_m(D)(1-t\mu)^{8-m}
\]
and using \eqref{eq:em-bounds}, with half-integral powers rounded down by
$q_S^{5/2}\le q_S^2$, $q_S^{3/2}\le q_S$, and
$q_S^{1/2}\le1$, gives
\begin{equation}\label{eq:IJ-middle}
 |I-J|\le E_S(a,\eta).
\end{equation}

Since
\[
 J=\frac{1-(1-a\mu)^9}{9\mu},
\]
write $R=|1-a\mu|$.  The bound $u\ge L_S(a)$ gives
\[
 R^2\le h_S(a,\eta).
\]
For a feasible pair $(a,\eta)$, this implies $h_S\ge0$, while
\eqref{eq:interval-hyp} gives $h_S\le1$.  Therefore
\[
 |J|\ge\frac{1-R^9}{9|\mu|}
 \ge\frac1{9\sqrt{1-\eta^2}}-
 \frac{h_S(a,\eta)^4}{9\lambda}
 \ge\frac19+\frac{\eta^2}{18}-
 \frac{h_S(a,\eta)^4}{9\lambda}.
\]
Together with \eqref{eq:IJ-middle}, this yields
\[
 |I|-\frac a9\ge P_{S,\lambda}(a,\eta)>0,
\]
contrary to \eqref{eq:I-upper}.
\end{proof}

We now give the exact finite certificate.  In each row of
Table~\ref{tab:middle}, Lemma~\ref{lem:sigma}, with $M=1+\beta$, gives
$\sigma\le S$.  The inequality $L_S(a)\ge\lambda$ is equivalent to
\[
 Sa^2-8\lambda a+8-S\ge0,
\]
which is checked exactly on the indicated interval.  The remaining elementary
conditions in \eqref{eq:interval-hyp} are direct rational inequalities.

\begin{table}[H]
\centering
\small
\renewcommand{\arraystretch}{1.18}
\caption{Exact constants for the middle range.}
\label{tab:middle}
\begin{tabular}{@{}c c c c@{}}
\toprule
$[\alpha,\beta]$ & $S$ & $\lambda$ & $Y$\\
\midrule
$[9/20,19/40]$ & $11043/2000$ & $49/50$ & $199/1000$\\
$[19/40,1/2]$ & $2783/500$ & $239/250$ & $587/2000$\\
$[1/2,21/40]$ & $2257/400$ & $1863/2000$ & $91/250$\\
$[21/40,11/20]$ & $1151/200$ & $1811/2000$ & $849/2000$\\
$[11/20,23/40]$ & $11819/2000$ & $879/1000$ & $477/1000$\\
$[23/40,3/5]$ & $764/125$ & $1699/2000$ & $66/125$\\
$[3/5,5/8]$ & $637/100$ & $817/1000$ & $577/1000$\\
$[5/8,13/20]$ & $13093/2000$ & $401/500$ & $239/400$\\
$[13/20,27/40]$ & $13113/2000$ & $81/100$ & $1173/2000$\\
$[27/40,7/10]$ & $3289/500$ & $409/500$ & $1151/2000$\\
$[7/10,29/40]$ & $529/80$ & $413/500$ & $141/250$\\
$[29/40,3/4]$ & $6661/1000$ & $1669/2000$ & $1103/2000$\\
$[3/4,31/40]$ & $1681/250$ & $843/1000$ & $269/500$\\
$[31/40,4/5]$ & $1701/250$ & $213/250$ & $131/250$\\
$[4/5,33/40]$ & $2761/400$ & $1723/2000$ & $127/250$\\
$[33/40,17/20]$ & $14041/2000$ & $109/125$ & $49/100$\\
$[17/20,7/8]$ & $7161/1000$ & $221/250$ & $187/400$\\
$[7/8,9/10]$ & $3663/500$ & $359/400$ & $883/2000$\\
\bottomrule
\end{tabular}
\end{table}

\begin{lemma}[middle-range certificate]\label{lem:middle-cert}
For every row $(\alpha,\beta,S,\lambda,Y)$ of
Table~\ref{tab:middle},
\[
 P_{S,\lambda}(a,\eta)>\frac1{2100}
 \qquad
 (\alpha\le a\le\beta,\ 0\le\eta\le Y).
\]
\end{lemma}

\begin{proof}
Substitute
\[
 a=\alpha+(\beta-\alpha)U,
 \qquad
 \eta=YV,
 \qquad 0\le U,V\le1.
\]
The resulting polynomial has bidegree at most $(9,8)$.  Its exact
Bernstein expansion is
\[
 \sum_{i=0}^{9}\sum_{j=0}^{8}b_{ij}
 \binom9iU^i(1-U)^{9-i}
 \binom8jV^j(1-V)^{8-j}.
\]
Across all eighteen rows and all $1620$ coefficients, the least coefficient is
\[
 \frac{1525995182936088798318336809}
 {3199338912963867187500000000000}
 >\frac1{2100}.
\]
The Bernstein basis functions are nonnegative and sum to $1$, proving the
claim.  The expansion and all preceding interval inequalities are verified in
exact rational arithmetic by the ancillary file.
\end{proof}

Proposition~\ref{prop:interval} and Lemma~\ref{lem:middle-cert} exclude the
entire interval $9/20\le a\le9/10$.

\section{The boundary range \texorpdfstring{$9/10\le a\le1$}{9/10 <= a <= 1}}

At $a=1$, \eqref{eq:localization} gives $u\ge1$, contradicting
$|\mu|<1$.  Suppose henceforth that $9/10\le a<1$, and put
\[
 \delta=1-a.
\]
Applying Lemma~\ref{lem:sigma} with $M=2$ gives
\begin{equation}\label{eq:end-sigma}
 \sigma\le\frac3{m_0^2}+1+\frac{256m_0^6}{81}<\frac{39}{5}.
\end{equation}
It follows from \eqref{eq:localization} that
\begin{equation}\label{eq:end-u}
 1-u\le\delta\frac{39a-1}{40a}\le\frac{19}{20}\delta.
\end{equation}
Consequently
\begin{equation}\label{eq:end-eta}
 0<\eta^2=1-|\mu|^2
 \le\frac{19}{10}\delta-\frac{361}{400}\delta^2.
\end{equation}

Write
\[
 \delta=x^2,
 \qquad
 \eta=xy.
\]
Then
\begin{equation}\label{eq:end-box}
 0<x<\frac8{25},
 \qquad
 0\le y<\frac75.
\end{equation}
For $0\le t\le1-x^2$, set
\begin{equation}\label{eq:end-q}
 q(x,y;t)=(1-t)^2+x^2t\left(\frac{19}{10}-y^2t\right).
\end{equation}
By \eqref{eq:end-u},
$|1-t\mu|^2\le q(x,y;t)$, and the same elementary calculation as before gives
$0\le q\le1$ on the feasible set.  Define
\begin{equation}\label{eq:end-E}
 \mathcal E(x,y)=
 \sum_{m=2}^{8}c_mx^{m-2}y^m
 \int_0^{1-x^2}t^m q(x,y;t)^{\ell_m}\,dt.
\end{equation}
The centered expansion and \eqref{eq:em-bounds} yield
\begin{equation}\label{eq:end-IJ}
 |I-J|\le\delta\mathcal E(x,y),
 \qquad
 J=\int_0^a(1-t\mu)^8\,dt.
\end{equation}

Let $R=|1-a\mu|$.  From \eqref{eq:end-u},
\[
 R^2\le
 \delta^2+\frac{19}{10}(1-\delta)\delta-(1-\delta)^2\eta^2
 \le\frac{19}{10}\delta.
\]
Also $|\mu|\ge181/200$.  Hence
\[
 |J|
 \ge\frac19+\frac{\eta^2}{18}-\frac{R^9}{9|\mu|}
 \ge\frac19+\frac{\eta^2}{18}-\frac{\delta}{1000}.
\]
For the last inequality, it is enough, by monotonicity in
$0<\delta\le1/10$, to check
\[
 \frac{200}{1629}\left(\frac{19}{10}\right)^{9/2}
 \left(\frac1{10}\right)^{7/2}\le\frac1{1000};
\]
after squaring this is the exact integer inequality
\[
 200^2 19^9<1629^2 10^{10}.
\]
Combining this with \eqref{eq:end-IJ} gives
\begin{equation}\label{eq:end-final-reduction}
 |I|-\frac a9
 \ge\delta\,\mathcal P(x,y),
\end{equation}
where
\begin{equation}\label{eq:end-P}
 \mathcal P(x,y)=
 \frac19+\frac{y^2}{18}-\frac1{1000}-\mathcal E(x,y).
\end{equation}
This is a rational polynomial of bidegree $(24,8)$.

\begin{lemma}[boundary certificate]\label{lem:end-cert}
For
\[
 0\le x\le\frac8{25},
 \qquad
 0\le y\le\frac75,
\]
one has
\[
 \mathcal P(x,y)>\frac1{100}.
\]
\end{lemma}

\begin{proof}
Put $x=(8/25)U$ and $y=(7/5)V$.  In the tensor-product Bernstein basis of
bidegree $(24,8)$, all $225$ coefficients are at least
\[
 \frac{139392070260030830922829855080988292511}
 {11102230246251565404236316680908203125000}
 >\frac1{100}.
\]
As in Lemma~\ref{lem:middle-cert}, positivity follows from the partition of
unity property of the Bernstein basis.  The coefficients are computed exactly
in the ancillary verifier.
\end{proof}

Equations \eqref{eq:I-upper} and \eqref{eq:end-final-reduction}, together with
Lemma~\ref{lem:end-cert}, are contradictory.  Thus no counterexample exists
for $9/10\le a\le1$.

\section{Completion of the proof}

The ranges $0\le a\le9/20$, $9/20\le a\le9/10$, and
$9/10\le a\le1$ have all been excluded under the assumption that every
critical point is farther than $1$ from $a$.  This proves
Theorem~\ref{thm:main}.

\appendix
\section{Exact Bernstein verification}

For reference, if
\[
 R(U,V)=\sum_{k=0}^{N}\sum_{\ell=0}^{M}a_{k\ell}U^kV^\ell,
\]
then its tensor-product Bernstein coefficients of bidegree $(N,M)$ are
\[
 b_{ij}=
 \sum_{k\le i}\sum_{\ell\le j}
 a_{k\ell}
 \frac{\binom{i}{k}}{\binom{N}{k}}
 \frac{\binom{j}{\ell}}{\binom{M}{\ell}}.
\]
This identity gives an especially short exact certificate: once all $b_{ij}$
are positive, $R$ is positive on $[0,1]^2$.  The accompanying verifier builds
the integrals in \eqref{eq:ES} and \eqref{eq:end-E} symbolically, performs the
affine substitutions, applies the displayed conversion formula using rational
numbers, and checks the stated coefficient minima.  A second implementation
independently reconstructs each polynomial from its Bernstein expansion.

\begin{thebibliography}{9}

\bibitem{BrandenWagner}
P.~Br\"and\'en and D.~G.~Wagner,
\newblock A converse to the Grace--Walsh--Szeg\H{o} theorem,
\newblock \emph{Math. Proc. Cambridge Philos. Soc.} \textbf{147} (2009),
447--453.

\end{thebibliography}

\end{document}
